% ═══════════════════════════════════════════════════════════════
% LH-08a: Lehrerhinweise - Viajar -- Reisen und Verkehrsmittel
% Abgeleitet aus LE-08a (Verlauf, Differenzierung, Fehler)
% ═══════════════════════════════════════════════════════════════

\documentclass[11pt, a4paper]{article}

\usepackage[utf8]{inputenc}
\usepackage[T1]{fontenc}
\usepackage[ngerman]{babel}
\usepackage{helvet}
\renewcommand{\familydefault}{\sfdefault}

\usepackage[margin=2cm]{geometry}
\usepackage{enumitem}
\usepackage{tabularx}
\usepackage{booktabs}
\usepackage{longtable}

\usepackage{xcolor}
\usepackage{tcolorbox}
\usepackage{amssymb}

\usepackage{fancyhdr}
\usepackage{lastpage}

% Farben
\definecolor{spanisch}{HTML}{c41e3a}
\definecolor{grau}{HTML}{888888}
\definecolor{hellgrau}{HTML}{f5f5f5}
\definecolor{basis}{HTML}{2a7a4b}
\definecolor{standard}{HTML}{2c5aa0}
\definecolor{erweiterung}{HTML}{7b1fa2}
\definecolor{loesung}{HTML}{c62828}

\newcommand{\fachfarbe}{spanisch}

% Variablen
\newcommand{\thema}{Viajar -- Reisen und Verkehrsmittel}
\newcommand{\sequenz}{08}
\newcommand{\block}{a}
\newcommand{\fach}{Spanisch}
\newcommand{\klasse}{7}
\newcommand{\dauer}{90 Min. (Doppelstunde)}

% Kopf-/Fußzeile
\pagestyle{fancy}
\fancyhf{}
\renewcommand{\headrulewidth}{0.4pt}
\renewcommand{\footrulewidth}{0.4pt}
\lhead{\fach{} -- Klasse \klasse}
\chead{\textbf{LH-\sequenz\block{} (Lehrerhinweise)}}
\rhead{}
\lfoot{}
\rfoot{Seite \thepage\ von \pageref{LastPage}}

% Differenzierungsmarker
\newcommand{\B}{\textcolor{basis}{\textbf{[B]}}}
\newcommand{\Std}{\textcolor{standard}{\textbf{[S]}}}
\newcommand{\E}{\textcolor{erweiterung}{\textbf{[E]}}}
\newcommand{\loesung}[1]{\textcolor{loesung}{\textbf{#1}}}
\newcommand{\es}[1]{\textcolor{\fachfarbe}{\textbf{#1}}}

% Boxen
\newtcolorbox{kurzplan}{
    colback=hellgrau,
    colframe=\fachfarbe,
    boxrule=1pt,
    arc=0pt,
    fonttitle=\bfseries,
    title=Kurzplan
}

\newtcolorbox{differenzierung}{
    colback=white,
    colframe=grau,
    boxrule=0.5pt,
    arc=0pt,
    fonttitle=\bfseries,
    title=Differenzierung
}

\newtcolorbox{fehlerbox}{
    colback=white,
    colframe=loesung,
    boxrule=0.5pt,
    arc=0pt,
    fonttitle=\bfseries,
    title=Häufige Fehler
}

\newtcolorbox{materialbox}{
    colback=white,
    colframe=\fachfarbe,
    boxrule=0.5pt,
    arc=0pt,
    fonttitle=\bfseries,
    title=Material-Checkliste
}

% ═══════════════════════════════════════════════════════════════
\begin{document}

\begin{center}
    {\Large\bfseries\textcolor{\fachfarbe}{Lehrerhinweise: \thema}}\\[0.3em]
    {\normalsize LH-\sequenz\block{} | \dauer{} | Std. 1-2 von 8}
\end{center}

\vspace{10pt}

% ═══════════════════════════════════════════════════════════════
% 1. KURZPLAN (aus LE-08a Verlauf)
% ═══════════════════════════════════════════════════════════════

\section*{1. Stundenverlauf (Kurzplan)}

\textbf{1. Stunde (45 Min.)}

\begin{tabularx}{\textwidth}{|c|l|X|l|}
\hline
\textbf{Zeit} & \textbf{Phase} & \textbf{Inhalt} & \textbf{Material} \\
\hline
0--5 & Einstieg & Bildcollage: ``¿Qué veis?'' & Bildcollage \\
\hline
5--15 & Input 1 & Verkehrsmittel einführen & Bildkarten \\
\hline
15--25 & Übung 1 & Memory: Bild-Wort-Zuordnung & Memory-Karten \\
\hline
25--35 & Input 2 & Orte + Aktivitäten einführen & Tafel, AB \\
\hline
35--45 & Übung 2 & AB-08a Aufgabe 1+2 & AB-08a \\
\hline
\end{tabularx}

\vspace{1em}

\textbf{2. Stunde (45 Min.)}

\begin{tabularx}{\textwidth}{|c|l|X|l|}
\hline
\textbf{Zeit} & \textbf{Phase} & \textbf{Inhalt} & \textbf{Material} \\
\hline
0--5 & Aktivierung & Quiz: Verkehrsmittel benennen & Bildkarten \\
\hline
5--15 & Input 3 & ir a + Infinitiv wiederholen & Tafel, AB \\
\hline
15--25 & Übung 3 & AB-08a Aufgabe 3: Sätze bilden & AB-08a \\
\hline
25--35 & Anwendung & PA: Dialog über Urlaubspläne & Karten \\
\hline
35--40 & Erweiterung & \E{}: Reiseplan schreiben & AB-08a \\
\hline
40--45 & Sicherung & Merkkasten, HA & AB-08a \\
\hline
\end{tabularx}

% ═══════════════════════════════════════════════════════════════
% 2. DIFFERENZIERUNG (aus LE-08a)
% ═══════════════════════════════════════════════════════════════

\section*{2. Differenzierung}

\begin{differenzierung}
\textbf{Legende:} \B{} Basis (BR) | \Std{} Standard (MR) | \E{} Erweiterung (GY)

\vspace{0.5em}
\textbf{WICHTIG:} Diese Marker sind NUR für die Lehrkraft!
\end{differenzierung}

\vspace{1em}

\begin{tabularx}{\textwidth}{|l|X|X|X|}
\hline
& \B{} \textbf{Basis} & \Std{} \textbf{Standard} & \E{} \textbf{Erweiterung} \\
\hline
\textbf{Aufg. 1} & Mit Wortbank & Ohne Wortbank & + Artikel ergänzen \\
\hline
\textbf{Aufg. 2} & Mit Bildhilfe & Nur Text & + eigene Zuordnungen \\
\hline
\textbf{Aufg. 3} & Wortbank für Verben & Nur Infinitiv gegeben & Ohne Hilfe \\
\hline
\textbf{Aufg. 4} & Dialog mit Vorlage & Dialog mit Stichworten & Freier Dialog \\
\hline
\textbf{Aufg. 5} & 3 Sätze, mit Hilfe & 5 Sätze & 8+ Sätze, + Kultur \\
\hline
\end{tabularx}

\vspace{1em}

\textbf{Konkrete Hinweise:}
\begin{itemize}
    \item \B{} LRS-SuS: Zeitzugabe (+10 Min.), Vokabelliste als Hilfe erlauben
    \item \B{} DaZ-SuS: Deutsch-Spanisch Wortschatz vorab besprechen
    \item \Std{} Standardgruppe: Partnerarbeit mit wechselnden Partnern
    \item \E{} Leistungsstarke: Kulturvergleich (Reisen in Spanien vs. Lateinamerika)
\end{itemize}

% ═══════════════════════════════════════════════════════════════
% 3. HÄUFIGE FEHLER (aus LE-08a Abschnitt 7)
% ═══════════════════════════════════════════════════════════════

\section*{3. Häufige Fehler}

\begin{fehlerbox}
\begin{tabularx}{\textwidth}{|l|X|X|}
\hline
\textbf{Fehler} & \textbf{Beispiel} & \textbf{Reaktion} \\
\hline
Aussprache \textit{avión} & *[a-vi-on] statt [a-bjon] & Vorsprechen, Silben klatschen: a-vi-ón \\
\hline
Verwechslung voy/vas/va & *``Yo vas a...'' & Kontrastive Übung: Pronomen + ir \\
\hline
Fehlende Präposition \textit{en} & *``Voy avión'' & Regel: en + Verkehrsmittel \\
\hline
Infinitiv ohne \textit{a} & *``Voy nadar'' & Merkspruch: ``ir A + Verb'' \\
\hline
Artikel bei Ländern & *``Voy a la España'' & Regel: Länder meist ohne Artikel \\
\hline
\end{tabularx}
\end{fehlerbox}

% ═══════════════════════════════════════════════════════════════
% 4. MATERIAL-CHECKLISTE
% ═══════════════════════════════════════════════════════════════

\section*{4. Material-Checkliste}

\begin{materialbox}
\begin{itemize}[label=$\square$]
    \item AB-08a.pdf (24 Kopien)
    \item ML-08a.pdf (1× für Lehrkraft)
    \item Lehrwerk: Qué pasa 1, S. 124--128
    \item Bildcollage: Urlaubsfotos (Strand, Berge, Flugzeug...)
    \item Bildkarten: Verkehrsmittel (12 Stück)
    \item Memory-Karten: Bild + Wort (12 Sets à 10 Karten)
    \item Dialogkarten mit Redemitteln
    \item Optional: LP-08a.html (Tablets/PCs)
\end{itemize}
\end{materialbox}

% ═══════════════════════════════════════════════════════════════
% 5. LÖSUNGEN (aus LE-08a Abschnitt 9)
% ═══════════════════════════════════════════════════════════════

\section*{5. Lösungen AB-08a}

\textbf{Merkkasten 1 (Verkehrsmittel):}
\begin{itemize}
    \item el avión, el tren, el coche, el autobús, el barco
    \item Präposition: \textbf{en} (Vamos \textbf{en} avión.)
\end{itemize}

\textbf{Merkkasten 2 (ir a + Infinitiv):}
\begin{itemize}
    \item yo $\rightarrow$ voy, tú $\rightarrow$ vas, él/ella $\rightarrow$ va, nosotros $\rightarrow$ vamos
\end{itemize}

\vspace{0.5em}

\textbf{Aufgabe 1:} (je 1P)
\begin{enumerate}[label=\alph*)]
    \item \loesung{el avión}
    \item \loesung{el tren}
    \item \loesung{el coche}
    \item \loesung{el autobús}
    \item \loesung{el barco}
\end{enumerate}

\textbf{Aufgabe 2:} (je 1P)
\begin{itemize}
    \item la playa $\rightarrow$ nadar, tomar el sol
    \item la montaña $\rightarrow$ hacer senderismo
    \item el hotel $\rightarrow$ descansar
    \item el camping $\rightarrow$ hacer una barbacoa
\end{itemize}

\textbf{Aufgabe 3:} (je 1P)
\begin{enumerate}[label=\alph*)]
    \item \loesung{voy a nadar}
    \item \loesung{vamos a viajar}
    \item \loesung{vas a hacer}
    \item \loesung{van a tomar}
    \item \loesung{va a ir}
    \item \loesung{vais a quedaros}
\end{enumerate}

\textbf{Aufgabe 4:} (6P)
\begin{itemize}
    \item Ziel nennen: Voy a ir a la playa / a España. (2P)
    \item Verkehrsmittel: Voy a viajar en avión / en coche. (2P)
    \item Aktivität: Voy a nadar / tomar el sol. (2P)
\end{itemize}

\textbf{Aufgabe 5:} (3P)
\begin{itemize}
    \item Ziel: 0,5P
    \item Verkehrsmittel: 0,5P
    \item Unterkunft: 0,5P
    \item Aktivitäten: 1P
    \item Korrektheit: 0,5P
\end{itemize}

% ═══════════════════════════════════════════════════════════════
% 6. HAUSAUFGABE
% ═══════════════════════════════════════════════════════════════

\section*{6. Hausaufgabe}

\begin{itemize}
    \item Lehrwerk S. 126, Übung 4 (Vokabeln)
    \item Vokabeln lernen: Verkehrsmittel, Orte, Aktivitäten
    \item Optional: Lernpfad LP-08a.html bearbeiten
    \item \E{}: Reiseplan erweitern (10 Sätze)
\end{itemize}

% ═══════════════════════════════════════════════════════════════
% 7. REFLEXIONSNOTIZEN
% ═══════════════════════════════════════════════════════════════

\section*{7. Reflexionsnotizen (nach Durchführung)}

\begin{tabularx}{\textwidth}{|l|X|}
\hline
\textbf{Was lief gut?} & \\[2em]
\hline
\textbf{Was lief nicht?} & \\[2em]
\hline
\textbf{Zeitmanagement} & \\[2em]
\hline
\textbf{Für nächstes Mal} & \\[2em]
\hline
\end{tabularx}

\end{document}
